\section{Other geometries} \label{secother}
The Segre integrals for arbitrary surfaces are established only when $F$ has rank $1$ or $2$, see \cite{MOP1, MOP}. In rank $1$, the answers were conjectured by Lehn \cite{Le}. The formulation below can be found in \cite{MOP}:
\begin{equation}\label{lehnc}
\sum_{n=0}^{\infty} z^k \int_{X^{[k]}} s\big(L^{[k]}\big) = A_1 (z)^{L^2} A_2(z)^{\chi(\mathcal O_X)} A_3 (z)^{L.K_X} A_4(z)^{K_X^2}.
\end{equation}
Here, for $z=t(1+2t)^2$, we have
\begin{gather*}
A_1(z)=(1+2t)^{\frac{1}{2}},
\\
A_2(z)=(1+2t)^{\frac{3}{2}} (1+6t)^{-\frac{1}{2}},
\\
A_3(z)=\frac{1}{2} (1+2t)^{-1} \left(\sqrt{1+2t}+\sqrt{1+6t}\right),
\\
 A_4(z)=4 (1+2t)^{\frac{1}{2}} (1+6t)^{\frac{1}{2}} \left(\sqrt{1+2t}+\sqrt{1+6t}\right)^{-2}.
 \end{gather*}

\subsection{A positivity result} It is more difficult to determine the sign of the top Segre integral from the above formulas. The~following lemma plays an important role in the analysis. We will apply it in the next section to geometric situations.

\begin{Lemma}\label{cl}
For all integers $m$, $n$, $p$ with $m\geq 0$, $p\geq 0$ such that $m+n+p$ is even, the series
\begin{gather*}
f(t)=\left(\sqrt{1+2t}+\sqrt{1+6t}\right)^m {(1+2t)}^{\frac{n-1}{2}} (1+6t)^{\frac{p-1}{2}}
\end{gather*}
has positive coefficients up to order less or equal than $\min\big(\frac{1}{2}(m+n+p)-1, m-1\big)$.
\end{Lemma}
Taking the minimum is necessary. Indeed, for $(m, n, p)=(2, 19, 1)$ the term of order $\frac{1}{2}(m+n+p)-1$ has negative coefficient, while for $(m, n, p)=(4, 0, 0)$, the term of order $m-1$ has zero coefficient. The lemma also holds for $m+n+p$ odd, but this case {\it never} occurs geometrically. The hypothesis $m, p\geq 0$ can fail in geometric examples, but our proof requires it.

\begin{proof}
The argument is not straightforward due to the alternating signs
in the expansions
\begin{gather*}
\sqrt{1+2t}= 1+t - \frac{t^2}{2}+\frac{t^3}{2}-\frac{5t^4}{8}+\cdots,
\\
\sqrt{1+6t}=1 + 3t - \frac{9t^2}{2}+\frac{27t^3}{2} - \frac{405 t^4}{8} +\cdots.
\end{gather*}
However, a residue calculation and suitable changes of variables will render the answer manifestly positive.

We begin by considering the case $n\geq 0$. In this situation, assume first $m$, $n$, $p$ are even, and~set
\begin{gather*}
F(t)=\left(\sqrt{1+2t}+\sqrt{1+6t}\right)^m {(1+2t)}^{-\frac{1}{2}} (1+6t)^{-\frac{1}{2}}.
\end{gather*}

\begin{Claim}\label{claim}
The series $F$ has positive coefficients up to order less or equal than $\frac{m}{2}-1$, and no terms of order between $\frac{m}{2}$ and $m-1$.
\end{Claim}


We assume this for now. Note that
\begin{gather*}
f(t)=F(t) P(t), \quad P(t)=(1+2t)^{\frac{n}{2}} (1+6t)^{\frac{p}{2}}.
\end{gather*}
Clearly, $P$ is a polynomial of degree $\frac{n+p}{2}$ with positive coefficients. This observation together with Claim~\ref{claim} gives the argument. Indeed, let $k\leq \min\big(\frac{1}{2}(m+n+p)-1, m-1\big)$. Writing $F_i$ and $P_j$ for the coefficients of $F$ and $P$, we have
\begin{gather*}
\operatorname{Coeff}_{ t^k} f(t)=\sum F_iP_j, \qquad\text{the sum ranging over}\quad i+j=k, \quad i\geq 0,\quad 0\leq j\leq \frac{n+p}{2}.
\end{gather*}
The sum is nonnegative since $i\leq k\leq m-1$, so $F_i\geq 0$ by the Claim~\ref{claim}, and $P_j> 0$. The sum is in fact strictly positive. Indeed, for $k\geq \frac{n+p}{2}$, it contains the term
\begin{gather*}
F_{k-\frac{n+p}{2}} P_{\frac{n+p}{2}}>0.
\end{gather*}
The last statement is also a consequence of Claim~\ref{claim}, using that $k-\frac{n+p}{2}\leq\frac{m}{2}-1$. When $k<\frac{n+p}{2}$, the sum contains the term $F_0 P_k=2^m P_k>0$.

\begin{proof}[Proof of Claim~\ref{claim}]
We seek to show that for all $k\leq \frac{m}{2}-1$, the residue
\begin{gather*}
\operatorname{Res}_{t=0} \frac{F(t)}{t^{k+1}} \,{\mathrm dt}>0.
\end{gather*}
The peculiar change of variables
\begin{gather*}
t=\frac{s(s+1)}{2\sqrt{3}s+(2+\sqrt{3})}
\end{gather*}
will simplify the calculation.
For convenience, set
\begin{gather*}
a=\frac{2-\sqrt{3}}{2\sqrt{3}}>0\qquad\text{so that}\quad t=\frac{s(s+1)}{2\sqrt{3}(s+a+1)}.
\end{gather*}
We note the following identities
\begin{gather*}
\sqrt{1+2t}=\frac{s+\frac{\sqrt{3}+1}{2}}{3^{\frac{1}{4}} (s+a+1)^{\frac{1}{2}}}, \qquad \sqrt{1+6t}=\frac{3^{\frac{1}{4}}\big(s+\frac{\sqrt{3}+1}{2\sqrt{3}}\big)}{(s+a+1)^{\frac{1}{2}}},
\\
\sqrt{1+2t}+\sqrt{1+6t}=\frac{(\sqrt{3}+1)(s+1)}{3^{\frac{1}{4}}(s+a+1)^{\frac{1}{2}}},\qquad
{\mathrm dt}=\frac{\big(s+\frac{\sqrt{3}+1}{2}\big)\big(s+\frac{\sqrt{3}+1}{2\sqrt{3}}\big)}{2\sqrt{3} (s+a+1)^2}\,{\mathrm ds}.
\end{gather*}
From here, we obtain
\begin{gather*}
\operatorname{Res}_{t=0} \frac{F(t)}{t^{k+1}} \,{\mathrm dt}=\operatorname{Res}_{s=0} \frac{G(s)}{s^{k+1}} \,{\mathrm ds},
\end{gather*}
where
\begin{gather*}
G(s)=\alpha (s+1)^{m-k-1} (s+a+1)^{k-\frac{m}{2}}, \qquad
\alpha=2^k \big(\sqrt{3}+1\big)^{m} \sqrt{3}^{k-\frac{m}{2}}>0.
\end{gather*}
A second change of variables will be needed next (this change of variables could have been carried out simultaneously with the first, but this is a price worth paying for readability). We~set
\begin{gather*}
s=\frac{(a+1)^2}{a} \frac{u}{1-\frac{a+1}{a}u}.
\end{gather*}
The reader can verify that
\begin{gather*}
s+1=\frac{1+(a+1)u}{1-\frac{a+1}{a}u}, \qquad s+a+1=\frac{a+1}{1-\frac{a+1}{a}u}, \qquad
{\mathrm ds}=\frac{(a+1)^2}{a} \frac{1}{\big(1-\frac{a+1}{a}u\big)^2}\,{\mathrm du}.
\end{gather*}
By direct calculation, we find
\begin{gather*}
\operatorname{Res}_{s=0} \frac{G(s)}{s^{k+1}}\,{\mathrm ds}=\operatorname{Res}_{u=0} \frac{H(u)}{u^{k+1}}\,{\mathrm du}
\end{gather*}
for
\begin{gather*}
H(u)=\beta (1+(a+1)u)^{m-k-1} \biggl(1-\frac{a+1}{a}u\biggr)^{k-\frac{m}{2}}, \qquad \beta=\alpha a^{k} (a+1)^{-k-\frac{m}{2}}>0.
\end{gather*}
The first term is a polynomial with positive coefficients $a_i$ given by binomial numbers, for $0\leq i\leq m-k-1$. The second term also has positive coefficients
\begin{gather*}
b_j=\biggl(-\frac{a+1}{a}\biggr)^j\binom{k-\frac{m}{2}}{j}>0
\end{gather*}
since $k-\frac{m}{2}<0$. Thus
\begin{gather*}
\operatorname{Coeff}_{u^k} H(u)=\!\sum a_i b_j,\quad \text{the sum ranging over}\!\quad i+j=k, \quad 0\leq i\leq m-k-1, \quad j\geq 0.
\end{gather*}
This coefficient is positive since $a_i>0$, $b_j> 0$ and the sum is non-empty (the term $a_0b_k=b_k>0$ appears in the sum).

When $m$ is even, and $\frac{m}{2}\leq k\leq m-1$, the expression $H(u)$ is a polynomial in $u$ of degree $\frac{m}{2}-1<k$. Hence, the coefficient of $u^k$ in $H(u)$ vanishes.
This proves the claim (and completes the argument when $m$, $n$, $p$ are even and $n\geq 0$).
\end{proof}

When $n\geq 0$, there are three other cases to consider, which require different choices for $F$ and $P$:
\begin{itemize}\itemsep=0pt
\item [$(i)$] when $m$ even, $n$, $p$ odd, we set
\begin{gather*}
F(t)=\left(\sqrt{1+2t}+\sqrt{1+6t}\right)^m, \qquad
P(t)=(1+2t)^{\frac{n-1}{2}} (1+6t)^{\frac{p-1}{2}},
\end{gather*}
\item [$(ii)$] when $m$ odd, $n$ even, $p$ odd, we set
\begin{gather*}
F(t)=\left(\sqrt{1+2t}+\sqrt{1+6t}\right)^m {(1+2t)}^{-\frac{1}{2}}, \qquad
P(t)=(1+2t)^{\frac{n}{2}} (1+6t)^{\frac{p-1}{2}},
\end{gather*}
\item [$(iii)$] when $m$ odd, $n$ odd, $p$ even, we set
\begin{gather*}
F(t)=\left(\sqrt{1+2t}+\sqrt{1+6t}\right)^m (1+6t)^{-\frac{1}{2}}, \qquad
P(t)=(1+2t)^{\frac{n-1}{2}} (1+6t)^{\frac{p}{2}}.
\end{gather*}
\end{itemize}
In case $(i)$, $F$ has positive coefficients up to order $\frac{m}{2}$, and no terms of order between $\frac{m}{2}+1$ and $m-1$.
In cases $(ii)$ and~$(iii)$, $F$ has positive coefficients up to order $\frac{m-1}{2}$, and no terms of order between $\frac{m+1}{2}$ and $m-1$. The arguments are similar, and we leave the details to the reader.

When $n<0$, the reasoning used to prove Claim~\ref{claim} also works to deal directly with the function
\begin{gather*}
f(t)=\left(\sqrt{1+2t}+\sqrt{1+6t}\right)^m {(1+2t)}^{\frac{n-1}{2}} (1+6t)^{\frac{p-1}{2}}.
\end{gather*}
Following exactly the same steps, first changing from $t$ to $s$ and then from $s$ to $u$, we obtain
\begin{align*}
H(u)={}&\gamma (1+(a+1)u)^{m-k-1} \biggl(1-\frac{a+1}{a}u\biggr)^{k-\frac{m+n+p}{2}} \biggl(1-\frac{(\sqrt{3}+1)^3}{4\sqrt{3}}u\biggr)^{n}
\\
&\times \biggl(1+\frac{\big(\sqrt{3}+1\big)^3}{4}u\biggr)^p,
\end{align*}
for $\gamma>0$. By the same arguments, this has positive coefficients when the exponents
\begin{gather*}
m-k-1\geq 0, \qquad k-\frac{m+n+p}{2}<0, \qquad n<0, \qquad p\geq 0,
\end{gather*}
which we assumed.
\end{proof}


